%قبل از شروع مسئله ترجیحاً چیزی ننویسید. اما اگر خواستید بنویسید، در قالب لیستی (فقرات و شمارش) بنویسید.
%برای اضافه کردن نمره، بعدی از شروع‌مسئله، [x] را اضافه کنید و به جای x نمره مورد نظر را قرار دهید. برای زیرمسئله‌ها هم به همین صورت است.
%\شروع‌مسئله[10]
\شروع‌مسئله [16]
%\متن‌سیاه{سوال عملی اول:}\\
%\متن‌سیاه{دراپ‌تست}\\
%\متن‌سیاه{ورودی:}\\
%در خط اول ورودی، ابتدا k و سپس m داده می‌شود.\\
%\متن‌سیاه{خروجی:}\\
%در تنها سطر خروجی، تعداد آیفون‌های مورد نیاز می‌آید. اگر در هیچ صورتی چنین کاری ممکن نبود ۱- چاپ می‌شود.\\
%\متن‌سیاه{محدودیت‌ها:}\\
%$1\leq k\leq 10^4$

%$m\leq10^9$
%\\
%\textbf{\متن‌سیاه{ورودی و خروجی نمونه}\\
%% تست اول
%\begin{latin}
% \begin{tabularx}{\textwidth}{|X | X|}
% \hline
% \textbf{stdin} &
% \textbf{stdout}\\
% \hline
% 100 100&
% 1\\
% \hline
% 
% \end{tabularx}
%\end{latin}
%
%% تست دوم
%
%\begin{latin}
% \begin{tabularx}{\textwidth}{|X | X|}
% \hline
% \textbf{stdin} &
% \textbf{stdout}\\
% \hline
% 340 20&
% 3\\
% \hline
% 
% \end{tabularx}
%\end{latin}}
% صورت اصلی مسئله را در این‌جا بنویسید.
در این مسئله می‌خواهیم مدل تعمیم یافتۀ الگوریتم مرتب‌سازی ادغامی را بررسی کنیم. در این الگوریتم به‌جای اینکه آرایۀ اولیه را به دو بخش تقسیم کنیم، آن را به $k$ بخش تقسیم می‌کنیم؛ هر کدام از بخش‌ها را به صورت بازگشتی مرتب کرده و سپس $k$ لیست مرتب شده را با هم ادغام می‌کنیم. برای ساده‌سازیِ تحلیل می‌توانید فرض کنید $n$ توان کاملی از $k$ است.
%اگر مسئله چند قسمت داشت به این صورت عمل کنید:
\شروع{زیرمسائل}
\زیرمسئله[4] % مسئلۀ ۱آ
با استفاده از داده‌ساختار هرم، الگوریتمی ارائه کنید که در زمان
$O(n \log{k})$
لیست‌ها را ادغام کرده و یک لیست مرتب $n$تایی تحویل دهد.
\زیرمسئله[8] % مسئلۀ ۱ب
حال فرض کنید الگوریتم شما فقط باید شامل تعدادی فراخوانی تابع
\lr{merge($a$, $b$)}
باشد که دو لیست $a$ و $b$ را ورودی گرفته و ادغام شدۀ آن‌ها را برمی‌گرداند. هزینۀ فراخوانی این تابع 
$\Theta(|a| + |b|)$
است که $|a|$ و $|b|$ به ترتیب اندازۀ لیست‌های $a$ و $b$ می‌باشند. مجدداً الگوریتمی ارائه دهید که در زمان
$O(n \log{k})$
لیست‌ها را ادغام کرده و یک لیست مرتب $n$ تایی تحویل دهد.
\زیرمسئله[4] % مسئلۀ ۱پ
فرض کنید از روش دوم برای ادغام $k$ لیست استفاده کنیم. رابطه بازگشتی زمان اجرای الگوریتم مرتب سازی ادغامی تعمیم یافته را بنویسید. سپس نشان دهید که زمان اجرای این الگوریتم (مستقل از مقدار $k$) حتما از
$\Omega(n \log{n})$
است.
%\زیرمسئله % مسئلۀ ۱ت
\پایان{زیرمسائل}

\solution{
	
آ) برای اینکار کافیست ابتدا از روی $k$ عضو اول لیست‌ها یک هرم کمینه بسازیم. راس‌های این هرم علاوه بر مقدار، یک اشاره‌گر به لیستی که از آن آمده‌اند نیز دارند. سپس در هر مرحله عضو کمینه هرم را استخراج کرده و به عنوان عضو بعدی لیست ادغام شده قرار می‌دهیم، سپس اگر عضوی از لیستی که عضو کمینه به آن تعلق داشت باقی‌مانده بود، عضو بعدی آن لیست را در هرم کمینه درج می‌کنیم. این کار را تا خالی شدن تمام لیست‌ها ادامه می‌دهیم. پیچیدگی زمانی ساخت این هرم 
$O(k)$
بوده و از آنجایی که هر عضو نیز دقیقا یک‌بار در هرم درج و حذف شده پیچیدگی زمانی کلی برابر است با
$O(k + n \log{k}) = O(n \log{k})$

ب) لیست‌ها را به 
$\lceil \frac{k}{2} \rceil$
گروه حداکثر دوتایی تقسیم می‌کنیم و اعضای این گروه‌ها را ادغام می‌کنیم. در مرحله بعدی 
$\lceil \frac{k}{2} \rceil$
لیست جدید داریم که مجدداً آن‌ها را دو به دو ادغام می‌کنیم و این کار را تا جایی ادامه می‌دهیم که تنها یک لیست باقی بماند. مراحل ادغام معادل با یک درخت دودویی پر هستند که برگ‌هایش لیست‌های اولیه می‌باشند. در هر عمقی از این درخت تعداد عملیات‌های انجام شده برای ادغام تمام لیست‌ها از 
$O(n)$
می‌باشد و با توجه به اینکه تعداد برگ‌ها $k$ است حداکثر ارتفاع این درخت
$\lceil \log{k} \rceil$
است. بنابراین پیچیدگی زمانی این الگوریتم نیز 
$O(n \log{k})$
است.

پ) به سادگی می‌توان نشان داد که پیچیدگی زمانی روش دوم از 
$\Theta(n \log{k})$
است. حال برای زمان اجرای مرتب‌سازی ادغامی تعمیم‌یافته داریم
$$T(n) = kT(\frac{n}{k}) + f(n)$$
و می‌دانیم
$$f(n) \in \Theta(n \log{k}) \rightarrow \exists C > 0 \ f(n) \geq Cn \log{k}$$
می‌خواهیم به کمک استقرا برروی $n$ ثابت کنیم 
$$\exists A > 0 \ T(n) \geq An \log{n}$$
بنابه فرض استقرا رابطه 
$$T(x) \geq Ax \log{x}$$
برای تمام $x$ های 
$k, k^2, k^3, ..., \frac{n}{k}$
برقرار است. آنگاه داریم
$$kT(\frac{n}{k}) \geq kA\frac{n}{k}\log{\frac{n}{k}} = An\log{n} - An\log{k}$$
$$T(n) \geq kT(\frac{n}{k}) + Cn\log{k} \geq An\log{n} - An\log{k} + Cn\log{k}$$
و اگر $A$ را به گونه‌ای انتخاب کنیم که
$A \leq C$
آنگاه خواهیم داشت.
$T(n) \geq An\log{n}$

همچنین پایه استقرا نیز حالت $n = k$ است که در این حالت نیز داریم
$$T(k) \geq kT(1) + Ck\log{k} \geq Ak\log{k}$$
بنابراین به کمک استقرا ثابت کردیم 
$$\exists A > 0 \ T(n) \geq An \log{n}$$
پس اثبات تمام است.

}




%%%%%%%%%%%%%%%%%%%%%%%%%%%%%%%%%%%%%%%%%%%%%%%%%%%%%
% برخی دستورات مفید برای تمرین‌ها. %
%%%%%%%%%%%%%%%%%%%%%%%%%%%%%%%%%%%%%%%%%%%%%%%%%%%%%

% فرمول‌های نمونه:
%$f(x) = \Theta(x)$
%$T(x, y) \leq \log(x) + 2^y + \binom{2n}{n}$
% {\tt code\_function}


% لیست موارد بدون شماره‌گذاری:
%\شروع{فقرات}
% \فقره مورد اول 
% \شروع{فقرات}
% \فقره مورد اول داخلی 
% \شروع{فقرات}
% \فقره اولین مورد خیلی داخلی 
% \فقره دومین مورد خیلی داخلی 
% \پایان{فقرات}
% \فقره مورد دوم داخلی 
% \پایان{فقرات}
% \فقره مورد دوم 
%\پایان{فقرات}

% لیست موارد شماره‌گذاری شده:
%\شروع{شمارش}
% \فقره یک مورد 
% \فقره یک مورد دیگر 
% \فقره باز هم یک مورد دیگر
%\پایان{شمارش}

% برای نوشتن یک فرمول چند خطی که یک شماره بخوره. برای حالت‌های دیگر به سایت زیر مراجعه کنید:
% https://www.overleaf.com/learn/latex/Aligning_equations_with_amsmath
%\begin{align} 
% \begin{split}
% (x+y)^3 &= (x+y)^2(x+y) \\
% &= (x^2+2xy+y^2)(x+y) \\
% &= (x^3+2x^2y+xy^2) + (x^2y+2xy^2+y^3) \\
% &= x^3+3x^2y+3xy^2+y^3
% \end{split} 
%\end{align}

% نوشتن ماتریس و آرایه دوبعدی
%\begin{align}
% A = 
% \begin{bmatrix}
% A_{11} & A_{21} \\
% A_{21} & A_{22}
% \end{bmatrix}
%\end{align}

% اضافه کردن تصویر با شرح. برای اضافه کردن تصویر بدون شرح، فقط دستور درج‌تصویر را استفاده کنید.
%\شروع{شکل}[h!]
%\تنظیم‌ازوسط
% \درج‌تصویر[پهنا=0.5\پهنای‌متن]{img.jpg}
%\شرح{شرح تصویر}
%\برچسب{تصویر:۱}
%\پایان{شکل}
%\پایان{زیرمسائل}

%کد زیر را ببینید.
%%\lstinputlisting{source_filename.py}
%\begin{latin}
%\begin{lstlisting}
%# Import the modules
%import sys
%import random
%
%ans = True
%
%while ans:
%question = raw_input("Ask the magic 8 ball a question: (press enter to quit) ")
%
%answers = random.randint(1,8)
%
%if question == "":
%sys.exit()
%
%elif answers == 1:
%print "It is certain"
%
%elif answers == 2:
%print "Outlook good"
%
%elif answers == 3:
%print "You may rely on it"
%
%elif answers == 4:
%print "Ask again later"
%
%elif answers == 5:
%print "Concentrate and ask again"
%
%elif answers == 6:
%print "Reply hazy, try again"
%
%elif answers == 7:
%print "My reply is no"
%
%elif answers == 8:
%print "My sources say no"
%\end{lstlisting}
%\end{latin}
